\section{Exporting results}
\label{sec:export}

One dialog writes everything RAFTcorr exports: data files, per-frame tables, pictures, animations
and an HTML report. Open it with \ui{Export Results} in the main window, with
\menu{File}{Export Results…} (\kbd{Ctrl+E}), or with \ui{Export Results} in the strain window, once
a run has results. Displacement can be exported before any strain is computed.

The window you open the dialog from sets the starting values of the picture settings:
\ui{Deformed frame} or \ui{Reference frame}, \ui{Show background image}, the hidden-background
fill, \ui{Hide untracked regions}, the opacity and the physical units. From the strain window, the
pictures also follow its \ui{Fill trimmed edges (display only)} (Section~\ref{sec:strain-viz});
from the main window, the trimmed band stays empty. Each field starts with the colormap and colour range
you last gave it in either window. The dialog stays open after an export, so you can write several
formats in turn; \ui{Close} closes it and stops an export still running.

\subsection{Output folder, file names and physical units}

\ui{OUTPUT FOLDER} takes the folder to write to: type it, or pick it with \ui{Browse…}, which also
remembers it for the next time. The export buttons stay disabled until a folder is set; the folder
is created if needed. \ui{Open Folder} opens it after the first successful export.

Every export writes new names, \file{<prefix>\_<kind>\_<timestamp>}: the prefix is the name of the
image folder (spaces and characters Windows does not allow become \file{\_}), the timestamp the
time of the click, \file{YYYYMMDDHHMMSS}. If a file exists already, the dialog lists it and asks
before replacing it.

\begin{center}
\small
\begin{tabularx}{\linewidth}{@{}l X@{}}
  \toprule
  Tab & Written in the output folder \\
  \midrule
  \ui{Data}      & \file{<p>\_parameters\_<t>.json} (always); \file{<p>\_results\_<t>.npz},
                   \file{.mat}, \file{.h5}; \file{<p>\_npz\_<t>/frame\_0001.npz} \ldots;
                   \file{<p>\_csv\_<t>/frame\_0001.csv} \ldots \\
  \ui{Images}    & \file{<p>\_images\_<t>/<field>/frame\_0001.png} \ldots, a folder per field \\
  \ui{Animation} & \file{<p>\_animation\_<t>/<field>.gif} or \file{.mp4}, a file per field \\
  \ui{Report}    & \file{<p>\_report\_<t>.html} \\
  \bottomrule
\end{tabularx}
\end{center}

\file{<field>} is the internal field name, for example \file{disp\_u} or \file{strain\_exx}; hover
over a field in the dialog to see it.

\begin{notebox}
File names, the \ui{FRAME RANGE} boxes and the report count frames from 1, like the frame slider
(\file{frame\_0001} is the reference image); the arrays inside the data files count from 0
(Section~\ref{sec:conventions}). Per-frame files are written for computed frames only, the
reference image included.
\end{notebox}

\ui{PHYSICAL UNITS} holds \ui{Enable physical units}, the \ui{Pixel size} with its unit (\file{nm},
\file{\textmu m}, \file{mm}, \file{cm}, \file{m}, \file{inch} \ui{/ pixel}) and the
\ui{Frame rate} in \ui{fps}, starting from the window's physical units
(Section~\ref{sec:physical-units}).
In pictures, animations and the report, displacements are scaled to the length unit and the
colorbar shows it. Data files stay in pixels: they record the calibration (\file{pixel\_size},
\file{length\_unit}, \file{frame\_rate}), and you multiply by \file{pixel\_size} to convert. Strain
is dimensionless and never scaled. With physical units off, the files record 1, \file{px} and no
frame rate.

\subsection{Data tab}

\screenshot{fig-export-data.png}{The \ui{Data} tab of the export dialog, below the output folder
and the physical units: the formats, the displacement and strain variables, and
\ui{Export Data}.}{fig:export-data}

\ui{FORMAT} chooses the files (Figure~\ref{fig:export-data}). \ui{NumPy Archive (.npz)} and
\ui{MATLAB v7.3 (.mat)} are on by default, \ui{HDF5 (.h5)} and \ui{CSV (per frame)} off. With NPZ
ticked, \ui{NPZ: one file per frame (default: single merged file)} writes one \file{.npz} per
computed frame into a folder of its own. \ui{DISPLACEMENT} and \ui{STRAIN} choose the variables,
with \ui{All} and \ui{None} buttons; \ui{U}, \ui{V}, \ui{Magnitude}, \ui{εxx}, \ui{εyy} and
\ui{εxy} are ticked by default. \ui{Residual} is available when the run recorded one, and the strain
group once you have run \ui{Compute Strain}. NPZ, MAT and HDF5 files always hold the grid
coordinates, the region of interest and the computed-frame flags as well, and the parameters file
is always written.

\subsubsection*{What the data files hold}

NPZ, MAT and HDF5 hold the same variables under the same names (Table~\ref{tab:export-variables}).
$T$ is the number of images, $h \times w$ the bounding box of the region of interest of frame~1
(the grid), and $h_s \times w_s$ the strain grid: the same box, taken every \ui{Output spacing}
pixels of the strain settings (Section~\ref{sec:strain}).

\begin{table}[htbp]
\centering
\small
\caption{Variables of the exported data files. Per-frame arrays are float32, coordinates float64,
flags uint8. The strain variables carry the names of the data files; the strain window shows them
as \ui{εxx}, \ui{εyy}, \ui{εxy}, \ui{ε₁}, \ui{ε₂}, \ui{ε₁ angle}, \ui{max shear}, \ui{von Mises},
\ui{ω rot} and \ui{ΔA/A₀}.}
\label{tab:export-variables}
\begin{tabularx}{\linewidth}{@{}l l l X@{}}
  \toprule
  Name & Shape & Unit & Content \\
  \midrule
  \file{U}, \file{V}        & $(T, h, w)$ & px & displacement, $+u$ right, $+v$ down \\
  \file{Magnitude}          & $(T, h, w)$ & px & $\sqrt{U^2 + V^2}$ \\
  \file{Residual}           & $(T, h, w)$ & grey levels & the residual (when ticked) \\
  \file{X}, \file{Y}        & $(h, w)$ & px & reference-image column and row of each point \\
  \file{ROI}                & $(h, w)$ & & region of interest of frame~1, 1 inside \\
  \file{Done}               & $(T)$ & & 1 where the frame was computed \\
  \file{exx}, \file{eyy}, \file{exy} & $(T, h_s, w_s)$ & & strain components (tensor shear) \\
  \file{e1}, \file{e2}      & $(T, h_s, w_s)$ & & principal strains \\
  \file{principal\_angle}, \file{rotation} & $(T, h_s, w_s)$ & deg & angle of \file{e1} from $+x$;
    rotation \\
  \file{max\_shear}, \file{von\_mises} & $(T, h_s, w_s)$ & & $(\varepsilon_1 - \varepsilon_2)/2$;
    von Mises equivalent strain \\
  \file{area\_change}       & $(T, h_s, w_s)$ & & area change $\Delta A/A_0$ \\
  \file{X\_strain}, \file{Y\_strain} & $(h_s, w_s)$ & px & coordinates of the strain points (with
    any strain variable) \\
  \file{frame\_names}       & $(T)$ & & image file names in frame order \\
  \bottomrule
\end{tabularx}
\end{table}

\begin{itemize}
  \item \textbf{NaN.} Displacement, magnitude and residual are NaN outside the region of interest,
    at points that could not be tracked (cut out by a later frame's region of interest, or gone out
    of the image) and in frames not computed. Strain is NaN outside the region of interest, where
    the virtual strain gauge found too few valid points, in the trimmed band along edges, and in
    frames without strain.
  \item \textbf{Frames and coordinates.} Index 0 is the reference image; its displacement is 0
    inside the region of interest. Every field is in reference (Lagrangian) coordinates. \file{X}
    and \file{Y} are integer pixel positions, the centre of the top-left pixel at $(0, 0)$.
  \item \textbf{Provenance.} Every file carries a JSON document: the RAFTcorr version, creation
    time, frame names, the reference of every frame (\file{ref\_indices}, 0-based), the run
    information (model name and SHA-256 of its file, iterations, correlation mode, device, the run
    settings as run, duration, notes), the grid, the strain settings, the calibration, the
    conventions, and the shape and unit of every variable.
\end{itemize}

In \textbf{HDF5}, each variable is a compressed dataset with a \file{units} attribute; the JSON
document is the root attribute \file{metadata}, beside simple attributes such as \file{n\_frames},
\file{pixel\_size} and \file{frame\_rate} (NaN when unknown). In \textbf{MATLAB}, \file{load} gives
\file{U} as a $T \times h \times w$ array; MATLAB counts from 1, so \file{U(n,:,:)} is frame $n$ of
the slider. The JSON is the text \file{metadata\_json} (read it with \file{jsondecode}), and
\file{frame\_names} has one line per image. In \textbf{NPZ}, the JSON is the string array
\file{metadata\_json}. A \textbf{per-frame NPZ} holds that frame's arrays without the time axis,
the coordinates and the JSON, whose key \file{frame} is the 0-based index.

\begin{lstlisting}[language=Python,showstringspaces=false]
import json, numpy as np
z = np.load("specimen_results_20261119103000.npz")
meta = json.loads(str(z["metadata_json"]))
u = z["U"][4] * meta["pixel_size"]     # frame 5 of the slider, in meta["length_unit"]
\end{lstlisting}

\begin{lstlisting}
S = load('specimen_results_20261119103000.mat');
meta = jsondecode(S.metadata_json);
u5 = squeeze(S.U(5,:,:));               % frame 5 of the slider, in pixels
\end{lstlisting}

The \textbf{parameters file} repeats the provenance as readable JSON. The \textbf{CSV tables} have
one row per point of the region of interest of frame~1: \file{x [px]}, \file{y [px]}, then the
chosen variables with their unit, \file{U [px]}, \ldots, \file{exx [-]}, \ldots,
\file{rotation [deg]}; an empty field is a NaN. With an \ui{Output spacing} above 1\,px the strain
points do not coincide with the displacement rows, so the strain columns go into a second file per
frame, \file{frame\_0001\_strain.csv}.

\begin{warnbox}
A CSV table has one row per pixel of the region of interest: a 1000\,$\times$\,1000\,px region gives
a million rows per frame. Use HDF5, MAT or NPZ for sequences, and CSV for a few frames.
\end{warnbox}

\subsection{Images tab}

\screenshot{fig-export-images.png}{The \ui{Images} tab: one row per field with its colormap, colour
range and opacity, then \ui{IMAGE SETTINGS} and \ui{FRAME RANGE}.}{fig:export-images}

The list at the top has one row per field (Figure~\ref{fig:export-images}): \ui{Export} (ticked by
default for \ui{U}, \ui{V}, \ui{εxx}, \ui{εyy}, \ui{εxy}), the field, its \ui{Colormap}, an
\ui{Auto} or \ui{Fixed} range with \ui{Min} and \ui{Max}, and the \ui{Opacity} ($\alpha$, 0
transparent to 1 opaque). Fields without data are greyed out. \ui{IMAGE SETTINGS}:

\begin{description}[leftmargin=0.6cm,style=nextline]
  \item[Format]
    \ui{PNG} (default), \ui{TIFF}, \ui{JPEG}, \ui{PDF} or \ui{SVG}. PDF and SVG keep the colorbar as
    vectors around an embedded picture of the field.
  \item[Resolution (long edge), JPEG quality, DPI]
    \ui{Full resolution} (default) writes the size set by \ui{DPI} (72 to 600, default 150; the
    field is drawn 7\,inches along its long edge, about 1050\,px at 150\,DPI). A cap of \ui{512 px}
    to \ui{2048 px} limits the long edge of each image and keeps the layout; for PDF and SVG it
    limits the embedded picture. \ui{JPEG quality}: 60 to 100, default 92.
  \item[Auto range]
    For \ui{Auto} rows: \ui{Each frame} (default) takes the range from each frame's own values, as
    the windows do; \ui{Whole sequence} takes one range from all exported frames, so a colour means
    the same value in every image. Use it, or a \ui{Fixed} range, for pictures shown side by side.
  \item[Include colorbar]
    On by default; styled on the \ui{Preview \& Colorbar} tab. It is labelled with the variable name
    of the data files and the unit, for example \file{u (mm)} or \file{exx (-)}.
  \item[Render as]
    \ui{Deformed frame} draws the field where each point is in that frame, over that frame's image;
    \ui{Reference frame} draws it at the positions in frame~1, over frame~1. Untick
    \ui{Show background image} to draw the field alone on the hidden-background fill. \ui{Crop to
    ROI} cuts the picture to the region of interest plus 20 image pixels.
  \item[Hide untracked regions]
    On: points without a result stay transparent. Off: in \ui{Reference frame} pictures they take
    the nearest value inside the region of interest, as on the canvas. Data files keep NaN either
    way.
\end{description}

\ui{FRAME RANGE} exports \ui{All frames}, or \ui{From} \ldots\ \ui{to}, numbered like the slider.
\ui{Export Images} writes one picture per computed frame and ticked field. When the images of the
run cannot be read, the pictures go on white and the \ui{LOG} says so.

\subsection{Animation tab}

The \ui{Animation} tab has its own field list and writes one animation per ticked field.
\ui{Format} is \ui{MP4} (H.264, the default) or \ui{GIF}. MP4 needs the \file{imageio} and
\file{imageio-ffmpeg} packages; without them the MP4 entry is greyed out. A GIF loops and has 255 colours
per frame. \ui{Resolution (long
edge)} defaults to \ui{1024 px}; \ui{Full resolution} draws the images at their own pixel size.
Keep a cap for GIF, whose size grows fast with resolution. \ui{FPS} is 1 to 60 (default 10).
\ui{Frame step} (1 to 20, default 1) keeps every $N$th frame and divides the written frame rate by
$N$, so the animation lasts as long, only choppier. \ui{Include colorbar}, \ui{Render as},
\ui{Crop to ROI} and \ui{Hide untracked regions} work as on the \ui{Images} tab. \ui{Auto} fields
get one colour range for the whole animation. GIF and MP4 have no transparency: a transparent
hidden background is written as white.

\subsection{Report tab}

\ui{Generate Report} writes one HTML file, pictures included, that opens in any browser without a
network connection. Under \ui{CONTENT}, all on by default: \ui{Parameter summary table} (the run and
strain parameters), \ui{Field statistics (min/max/mean/std per frame)} over the valid points of the
region of interest, and \ui{Sample field images} of every $N$th computed frame plus the last
(\ui{Sample every} 5 \ui{frames} by default, 1 to 100). The pictures use each field's row and the
settings of the \ui{Images} tab, with one colour range per field. \ui{FIELDS} chooses the fields;
\ui{U}, \ui{V}, \ui{Magnitude}, \ui{εxx}, \ui{εyy} and \ui{εxy} by default.

\subsection{Preview \& Colorbar tab}

This tab draws one frame with the code of the \ui{Images} export, at a lower resolution: what you
see is what the export writes. Pick a \ui{Field} ticked on the \ui{Images} tab and a \ui{Frame}
(from 1; it starts on frame~2). Transparency shows as a grey checkerboard. \ui{FIELD APPEARANCE}
(\ui{Colormap}, \ui{Range}, \ui{Min}, \ui{Max}, \ui{Opacity}) is the same setting as the field's
row on the \ui{Images} tab; \ui{Apply to all fields} copies the colormap, opacity and range mode,
not \ui{Min} and \ui{Max}, to every ticked field of the \ui{Images} and \ui{Animation} tabs.
\ui{COLORBAR STYLE} applies to every picture, animation and report image:

\begin{center}
\small
\begin{tabularx}{\linewidth}{@{}l l X@{}}
  \toprule
  Setting & Default & Choices \\
  \midrule
  \ui{Position}          & \ui{Right} & \ui{Right}, \ui{Left}, \ui{Top}, \ui{Bottom} \\
  \ui{Font size}, \ui{Font family} & 11\,pt, sans-serif & 6 to 32\,pt; sans-serif, serif,
    monospace \\
  \ui{Bar thickness}     & 0.03 & 0.02 to 0.25 of the picture's width (bar left or right) or
    height (top or bottom) \\
  \ui{Background}        & \ui{White} & \ui{Black} or \ui{White} behind the colorbar and around
    the picture; text takes the other colour \\
  \ui{Margin}, \ui{Margin color} & 0, \ui{White} & 0 to 0.30 of the long edge; \ui{White} or
    \ui{Black} \\
  \ui{Hidden background} & from the window & \ui{White}, \ui{Black} or \ui{Transparent} where the
    image would be; transparency is kept in PNG and TIFF, other formats get white \\
  \bottomrule
\end{tabularx}
\end{center}

\subsection{Running an export}

A running export shows a progress bar and the item being written; \ui{Cancel Export} stops it
after the current frame or file and keeps what was written. Data files are written to a
temporary file and moved into place when complete, so a cancelled or failed export never leaves a
half-written data file and never damages an existing one. Errors appear in red under the button
and in the \ui{LOG}.

\begin{tipbox}
The dialog exports the results it was opened with. After a new run, a new strain computation or
opening a project, it refuses to export and asks you to close it and open it again.
\end{tipbox}
